\section*{第 \ref{ch:g12:integ} 章 --- 积分}

\begin{solution}{exo:g12:integ:1}
\[
\int_1^2 \left(3x^2 - \frac{1}{x^2}\right)\dd x
= \left[x^3 + \frac1x\right]_1^2 = \left(8 + \tfrac12\right) - 2 = \frac{13}{2}.
\]
\[
\int_0^{\pi/2} \cos t\,\dd t = \bigl[\sin t\bigr]_0^{\pi/2} = 1 .
\]
设 $u = 2t+1$，$\frac{1}{2t+1} = \frac12\,\frac{u'}{u}$：
\[
\int_0^1 \frac{\dd t}{2t+1} = \frac12\bigl[\ln(2t+1)\bigr]_0^1
= \frac{\ln 3}{2}.
\]
\end{solution}

\begin{solution}{exo:g12:integ:2}
$x\,\eu^{x^2} = \frac12\,(x^2)'\,\eu^{x^2}$ 的原函数为
$\frac12\eu^{x^2} + c$。条件 $F(0) = 1$ 给出 $\frac12 + c = 1$，
于是
\[
F(x) = \frac{\eu^{x^2} + 1}{2}.
\]
\end{solution}

\begin{solution}{exo:g12:integ:3}
\[
\mu = \frac{1}{\pi}\int_0^\pi \sin t\,\dd t
= \frac{1}{\pi}\bigl[-\cos t\bigr]_0^\pi = \frac{2}{\pi} \approx 0.64 .
\]
从图形上看，底为 $\intcc{0}{\pi}$、高为
$\frac{2}{\pi}$ 的矩形与正弦曲线的一拱面积相等。
\end{solution}

\begin{solution}{exo:g12:integ:4}
取 $u' = 1$，$v = \ln t$，则 $u = t$，$v' = \frac1t$：
\[
\int_1^{\eu} \ln t\,\dd t = \bigl[t\ln t\bigr]_1^{\eu} - \int_1^{\eu} 1\,\dd t
= \eu - (\eu - 1) = 1 .
\]
取 $u' = \sin t$，$v = t$，则 $u = -\cos t$，$v' = 1$：
\[
\int_0^\pi t\sin t\,\dd t = \bigl[-t\cos t\bigr]_0^\pi + \int_0^\pi \cos t\,\dd t
= \pi + 0 = \pi .
\]
\end{solution}

\begin{solution}{exo:g12:integ:5}
两曲线的交点满足 $x^2 = x + 2$，\emph{即} $x \in \{-1, 2\}$，且
在 $\intcc{-1}{2}$ 上直线位于抛物线上方，面积为
\[
\int_{-1}^{2} \bigl(x + 2 - x^2\bigr)\dd x
= \left[\frac{x^2}{2} + 2x - \frac{x^3}{3}\right]_{-1}^{2}
= \frac{10}{3} - \left(-\frac{7}{6}\right) = \frac{9}{2}.
\]
\end{solution}

\begin{solution}{exo:g12:integ:6}
\emph{1.} $I = \bigl[\ln(1+t)\bigr]_0^1 = \ln 2$.

\emph{2.}
\[
I_n + I_{n+1} = \int_0^1 \frac{t^n(1 + t)}{1+t}\,\dd t
= \int_0^1 t^n \,\dd t = \frac{1}{n+1}.
\]
在 $\intcc{0}{1}$ 上，$0 \leq \dfrac{t^n}{1+t} \leq t^n$，由比较可得
$0 \leq I_n \leq \frac{1}{n+1}$。

\emph{3.} 由递推关系，经逐项相消的归纳可得
\[
I_n = (-1)^n\left(I_0 - \left(1 - \frac12 + \dots + \frac{(-1)^{n-1}}{n}\right)\right),
\]
\emph{即}
$\displaystyle\sum_{k=1}^{n} \frac{(-1)^{k-1}}{k} = \ln 2 - (-1)^n I_n$
（回顾 $I_0 = \ln 2$）。由 $\abs{I_n} \leq \frac{1}{n+1} \to 0$，
交错调和级数的和趋于 $\ln 2$。
\end{solution}

\begin{solution}{exo:g12:integ:7}
路程：
\[
d = \int_0^{100} 30\left(1 - \eu^{-t/50}\right)\dd t
= 30\left[t + 50\,\eu^{-t/50}\right]_0^{100}
= 30\left(100 + 50\eu^{-2} - 50\right)
= 1500 + 1500\,\eu^{-2}.
\]
数值上 $d \approx 1500 + 203 = 1703$ 米。平均速度：
$\frac{d}{100} = 15\bigl(1 + \eu^{-2}\bigr) \approx 17.0$ 米/秒。
\end{solution}

\begin{solution}{exo:g12:integ:8}
\emph{1.} $S_n$ 是 $n$ 个宽为 $\frac1n$ 的矩形面积之和，
高分别为 $f\left(\frac kn\right)$，$k = 1, \dots, n$，其中
$f(t) = \frac{1}{1+t}$：它是 “右端点”矩形法
$f$ 在 $\intcc{0}{1}$ 上曲线下方面积的近似。

\emph{2.} $f$ 递减，所以在每个区间
$\intcc{\frac{k-1}{n}}{\frac kn}$ 上，
$f\left(\frac kn\right) \leq f(t) \leq f\left(\frac{k-1}{n}\right)$；
对上述不等式积分并对 $k$ 求和：
\[
S_n \leq \int_0^1 f(t)\,\dd t \leq S_n + \frac1n\bigl(f(0) - f(1)\bigr)
= S_n + \frac{1}{2n}.
\]
（习题展示式中的中间项，
当 $f(0) - f(1) = 1 - \frac12$ 时，即为此界。）因此
$0 \leq \ln 2 - S_n \leq \frac{1}{2n} \to 0$，故 $S_n \to \ln 2$。
\end{solution}

\begin{solution}{exo:g12:integ:9}
\emph{1.} $W_0 = \int_0^{\pi/2} \dd t = \frac{\pi}{2}$ 且
$W_1 = \bigl[-\cos t\bigr]_0^{\pi/2} = 1$。

\emph{2.} 分部积分，取 $u' = \sin t$，$v = \sin^{n+1} t$，
从而 $u = -\cos t$，$v' = (n+1)\sin^n t \cos t$：
\[
W_{n+2} = \bigl[-\cos t \sin^{n+1} t\bigr]_0^{\pi/2}
+ (n+1)\int_0^{\pi/2} \cos^2 t\,\sin^n t\,\dd t
= (n+1)\int_0^{\pi/2} (1 - \sin^2 t)\sin^n t\,\dd t,
\]
两端的方括号项为零，于是
$W_{n+2} = (n+1)(W_n - W_{n+2})$，\emph{即}
$W_{n+2} = \frac{n+1}{n+2} W_n$。

\emph{3.} $W_2 = \frac12 W_0 = \frac{\pi}{4}$，
$W_3 = \frac23 W_1 = \frac23$，
$W_4 = \frac34 W_2 = \frac{3\pi}{16}$。
在 $\intoo{0}{\frac\pi2}$ 上，$0 < \sin t < 1$，故
$\sin^{n+1} t \leq \sin^n t$，在内部严格成立；
积分得 $0 < W_{n+1} \leq W_n$（事实上为 $<$）：$(W_n)$ 递减且
为正。
\end{solution}

\begin{solution}{pb:g12:integ:1}
\textbf{1.} $\left[x^3 - x^2 + x\right]_0^1 = 1$；\quad
$\left[\ln x\right]_1^{\eu} = 1$；\quad
$\left[\sin x\right]_0^{\pi/2} = 1$。

\textbf{2.} $\left[\frac12 \eu^{x^2}\right]_0^1 =
\frac{\eu - 1}{2}$；\quad
$\left[\ln(x^2 + 1)\right]_0^1 = \ln 2$。

\textbf{3.} $\int_0^1 x\eu^x = \left[x\eu^x\right]_0^1 -
\int_0^1 \eu^x = \eu - (\eu - 1) = 1$。同理有
$\int_1^{\eu} \ln x = \left[x\ln x\right]_1^{\eu} -
\int_1^{\eu} 1 = \eu - (\eu - 1) = 1$。

\textbf{4.} $\frac{1}{\pi}\int_0^\pi \sin x\,\dd x =
\frac{2}{\pi} \approx 0.64$：正弦拱的平均高度
为 $\frac2\pi$，明显大于一半——
拱在顶部附近较为饱满。

\textbf{5.} $\int_0^1 (x - x^2)\,\dd x = \frac12 - \frac13 =
\frac16$。

\textbf{6.} $L_n = \frac1n \sum_{k=0}^{n-1}
\left(\frac kn\right)^2$ 与
$U_n = \frac1n \sum_{k=1}^{n} \left(\frac kn\right)^2$：
宽为 $\frac1n$ 的矩形，高度取左端点
（位于曲线下方）或右端点（位于曲线上方）处的函数值。

\textbf{7.} 当 $n = 1$ 时成立（$1 = \frac{1 \cdot 2 \cdot
3}{6}$）。若对 $n$ 成立，则两边加上 $(n+1)^2$：
\[
\frac{n(n+1)(2n+1)}{6} + (n+1)^2
= \frac{(n+1)\left(2n^2 + n + 6n + 6\right)}{6}
= \frac{(n+1)(n+2)(2n+3)}{6},
\]
即为 $n + 1$ 时的公式，归纳步证毕。

\textbf{8.} $U_n = \frac{1}{n^3} \cdot
\frac{n(n+1)(2n+1)}{6} = \frac{(1 + \frac1n)(2 + \frac1n)}{6}
\to \frac26 = \frac13$，且 $L_n = U_n - \frac1n \to \frac13$：
夹在两组阶梯形之间，
面积恰为 $\frac13$。

\textbf{9.} $\int_0^1 x^2\,\dd x =
\left[\frac{x^3}{3}\right]_0^1 = \frac13$：只需一行。
微积分基本定理把无穷多次夹逼
化为一次原函数求值——这正是它被称为
“基本”的原因。

\textbf{10.} 抛物线弓形：$\int_{-1}^{1}(1 - x^2)\,\dd x =
2 - \frac23 = \frac43$。三角形以弦为底（长为
$2$，位于高 $1$ 处），顶点为 $(0,0)$，面积
为 $\frac12 \times 2 \times 1 = 1$。面积比为 $\frac43$——恰是
阿基米德定理，由他所缺的“机器”加以验证。

\textbf{11.} $T\left(1 + \frac14 + \frac{1}{16} + \dots\right)
= T \cdot \frac{1}{1 - \frac14} = \frac43 T$（等比级数
公比为 $\frac14$）：取 $T = 1$，抛物线弓形的面积
$\frac43$ 再次出现。阿基米德用纯几何论证
求和——那种“每次咬去四分之一”的
巧克力条式做法，比“极限”概念的形成早两千年。

\textbf{12.} 穷竭法需要巧妙的几何自相似性，
并逐个情形给出精确面积；黎曼
只需单调性，便能给出普遍的定义，
代价是计算和式；微积分基本定理则需要
导数的整个体系，并以此回报：使面积
只需一行即可计算。

\textbf{13.} $A(1) = 0$ 且 $A'(x) = \frac1x$（由
微积分基本定理）：$A$ 是 $\frac1x$ 在
$\intoo{0}{+\infty}$ 上在 $1$ 处为零的原函数——
恰为自然对数。

\textbf{14.} $g'(x) = \frac{a}{ax} - \frac1x = 0$：$g$ 为
常数；取 $x = 1$，$g(1) = A(a) - A(1) = A(a)$。于是
$A(ax) = A(a) + A(x)$ 对所有 $x$ 成立，取 $x = b$，
$\ln(ab) = \ln a + \ln b$。双曲线下的面积
与横向尺度无关：将 $t$ 伸长 $a$ 倍，
会使 $\frac1t$ 压缩 $a$ 倍，而面积不变——
对数的加法正是一种尺度对称。

\textbf{15.} 逐次迭代：$\ln(a^n) = n\ln a$（对
$n$ 用归纳法）。再由 $0 = \ln 1 = \ln\left(a \cdot \frac1a\right) =
\ln a + \ln\frac1a$：$\ln\frac1a = -\ln a$。

\textbf{16.} $\frac{1}{11} + \frac{1}{12} + \dots +
\frac{1}{20} \approx 0.669$，与 $\ln 2 \approx 0.693$ 相比：
十个位于 $\frac{1}{1+t}$ 下方的矩形已十分接近。调和级数
发散，但其从 $n$ 到 $2n$ 的\emph{截段}之和
趋于 $\ln 2$：这是借
面积来度量的一种缓慢发散。

\textbf{17.} 距离：$\int_0^{10} 3t^2\,\dd t = 1000$ m。
平均速率：$100$ m/s。瞬时速度为 $v = 100$ 的时刻：
$3t^2 = 100$：$t = \sqrt{\frac{100}{3}} \approx 5.77$ s。

\textbf{18.} 质量：$\int_0^4 (2 + x)\,\dd x = 8 + 8 = 16$ kg。
一阶矩：$\int_0^4 (2x + x^2)\,\dd x = 16 + \frac{64}{3} =
\frac{112}{3}$。质心：
$\frac{112/3}{16} = \frac73 \approx 2.33$ m --- 质心被较重的右端
推过中点，与直觉相符。

\textbf{19.} $W_0 = \frac\pi2$；$W_1 = \left[-\cos
t\right]_0^{\pi/2} = 1$；而
$W_2 = \int_0^{\pi/2} \frac{1 - \cos 2t}{2}\,\dd t =
\frac\pi4$。

\textbf{20.} 面积：由定义出发，用阶梯形夹逼。
累积：由速率求距离，由密度求质量 ---
积分即为逐步累加的总量。原函数：
微积分基本定理，牛顿与莱布尼茨借
阿基米德式耐心而得的捷径。点睛之笔：$\ln$ 键是
$\frac1t$ 下方的面积，而该面积的尺度对称\emph{正是}
恒等式 $\ln(ab) = \ln a + \ln b$。
\end{solution}
